
\documentclass{article}
\usepackage{colortbl}
\usepackage{makecell}
\usepackage{multirow}
\usepackage{supertabular}

\begin{document}

\newcounter{utterance}

\centering \large Interaction Transcript for game `codenames', experiment `association\_difficult', episode 4 with Apertus{-}v1.5{-}8B.
\vspace{24pt}

{ \footnotesize  \setcounter{utterance}{1}
\setlength{\tabcolsep}{0pt}
\begin{supertabular}{c@{$\;$}|p{.15\linewidth}@{}p{.15\linewidth}p{.15\linewidth}p{.15\linewidth}p{.15\linewidth}p{.15\linewidth}}
    \# & $\;$A & \multicolumn{4}{c}{Game Master} & $\;\:$B\\
    \hline

    \theutterance \stepcounter{utterance}  
    & & & \multicolumn{2}{p{0.3\linewidth}}{
        \cellcolor[rgb]{0.9,0.9,0.9}{
            \makecell[{{p{\linewidth}}}]{
                \texttt{\tiny{[GM$|$GM]}}
                {'hidden': {'team': ['цукор', 'ядерне паливо', 'лялька', 'квартира', 'корова', 'опера', 'лось', 'неетилований бензин', 'вігвам'], 'innocent': ['білка', 'нирка', 'кориця', 'блюграсс', 'масло', 'кит', 'кубики'], 'opponent': ['фолк', 'сонячна енергія', 'кумин', 'шкіра', 'сом', 'вітер', 'пазл', 'пума'], 'assassin': ['джакс']}, 'revealed': {'team': {'team': [], 'innocent': [], 'opponent': [], 'assassin': []}, 'opponent': {'team': [], 'innocent': [], 'opponent': [], 'assassin': []}}}
            }
        }
    }
    & & \\ \\

    \theutterance \stepcounter{utterance}  
    & & \multicolumn{4}{p{0.6\linewidth}}{
        \cellcolor[rgb]{0.9,0.9,0.9}{
            \makecell[{{p{\linewidth}}}]{
                \texttt{\tiny{[P1$\langle$GM]}}
                \texttt{Надайте одну однословну підказку, яка пов'язана з деякими словами вашої команди зі списку нижче. Ви можете обрати, зі скількома і з якими саме словами пов'язана ваша підказка.} \\
\texttt{Мета полягає в тому, щоб ваш партнер відгадав усі слова вашої команди на основі підказок, які ви йому даєте, раніше, ніж інша команда відкриє всі свої. Якщо ваш партнер відгадає слово «суперника» або «відволікача», його хід відгадування закінчується, і команда суперників отримує право зробити хід. Якщо ваш партнер відгадає слово «вбивці», гра закінчується, і ваша команда негайно програє.} \\
\texttt{Чим із більшою кількістю слів пов'язана ваша підказка, тим швидше ви можете виграти гру, але тим складнішими можуть бути асоціації для відгадування вашим партнером, тож обирайте слово{-}підказку мудро. Слово{-}підказка має бути семантично пов'язаним із цільовими словами; воно не може бути одним зі слів у списках або містити їхні частини.} \\
\\ 
\texttt{Завжди надавайте вашу однословну підказку та ваш розділений комами список пов'язаних цільових слів у такому форматі й робіть відповіді якомога коротшими; ніколи не додавайте жодного іншого тексту, окрім того, що потрібно в цій формі:} \\
\\ 
\texttt{ПІДКАЗКА: <СЛОВО>} \\
\texttt{ЦІЛІ: <ЦІЛІ>} \\
\\ 
\texttt{Слова вашої команди: цукор, ядерне паливо, лялька, квартира, корова, опера, лось, неетилований бензин, вігвам.} \\
\texttt{Слова команди суперника: фолк, сонячна енергія, кумин, шкіра, сом, вітер, пазл, пума.} \\
\texttt{Слова{-}відволікачі: білка, нирка, кориця, блюграсс, масло, кит, кубики.} \\
\texttt{Слова{-}вбивці: джакс.} \\
            }
        }
    }
    & \\ \\

    \theutterance \stepcounter{utterance}  
    & \multicolumn{4}{p{0.6\linewidth}}{
        \cellcolor[rgb]{0.9,0.9,0.9}{
            \makecell[{{p{\linewidth}}}]{
                \texttt{\tiny{[P1$\rangle$GM]}}
                \texttt{ПІДКАЗКА: енергетика} \\
\texttt{ЦІЛІ: ядерне паливо, сонячна енергія} \\
            }
        }
    }
    & & \\ \\

    \theutterance \stepcounter{utterance}  
    & & & \multicolumn{2}{p{0.3\linewidth}}{
        \cellcolor[rgb]{0.9,0.9,0.9}{
            \makecell[{{p{\linewidth}}}]{
                \texttt{\tiny{[GM$|$GM]}}
                \texttt{енергетика} \\
            }
        }
    }
    & & \\ \\

    \theutterance \stepcounter{utterance}  
    & & & \multicolumn{2}{p{0.3\linewidth}}{
        \cellcolor[rgb]{0.9,0.9,0.9}{
            \makecell[{{p{\linewidth}}}]{
                \texttt{\tiny{[GM$|$GM]}}
                ['ядерне паливо', 'сонячна енергія']
            }
        }
    }
    & & \\ \\

    \theutterance \stepcounter{utterance}  
    & & & \multicolumn{2}{p{0.3\linewidth}}{
        \cellcolor[rgb]{0.9,0.9,0.9}{
            \makecell[{{p{\linewidth}}}]{
                \texttt{\tiny{[GM$|$GM]}}
                {'word': 'ядерне паливо', 'assignment': 'team'}
            }
        }
    }
    & & \\ \\

    \theutterance \stepcounter{utterance}  
    & & & \multicolumn{2}{p{0.3\linewidth}}{
        \cellcolor[rgb]{0.9,0.9,0.9}{
            \makecell[{{p{\linewidth}}}]{
                \texttt{\tiny{[GM$|$GM]}}
                {'word': 'сонячна енергія', 'assignment': 'opponent'}
            }
        }
    }
    & & \\ \\

    \theutterance \stepcounter{utterance}  
    & & \multicolumn{4}{p{0.6\linewidth}}{
        \cellcolor[rgb]{0.9,0.9,0.9}{
            \makecell[{{p{\linewidth}}}]{
                \texttt{\tiny{[GM$\rangle$P2]}}
                \texttt{Надайте розділений комами список із не більш ніж 2 слів із наведеного нижче списку, які найкраще пов'язані або найбільш тісно асоціюються зі словом «енергетика». Завжди починайте свій список здогадок із ВІДГАДКА: і не додавайте жодного іншого тексту у своїй відповіді.} \\
\\ 
\texttt{нирка, квартира, лось, блюграсс, кориця, ядерне паливо, кит, білка, джакс, сом, цукор, фолк, пазл, опера, неетилований бензин, лялька, вітер, кумин, корова, вігвам, шкіра, масло, пума, кубики, сонячна енергія} \\
            }
        }
    }
    & \\ \\

    \theutterance \stepcounter{utterance}  
    & & & \multicolumn{4}{p{0.6\linewidth}}{
        \cellcolor[rgb]{0.9,0.9,0.9}{
            \makecell[{{p{\linewidth}}}]{
                \texttt{\tiny{[GM$\langle$P2]}}
                \texttt{ВІДГАДКА: ядерне паливо, сонячна енергія} \\
            }
        }
    }
     \\ \\

    \theutterance \stepcounter{utterance}  
    & & & \multicolumn{2}{p{0.3\linewidth}}{
        \cellcolor[rgb]{0.9,0.9,0.9}{
            \makecell[{{p{\linewidth}}}]{
                \texttt{\tiny{[GM$|$GM]}}
                ['ядерне паливо', 'сонячна енергія']
            }
        }
    }
    & & \\ \\

    \theutterance \stepcounter{utterance}  
    & & & \multicolumn{2}{p{0.3\linewidth}}{
        \cellcolor[rgb]{0.9,0.9,0.9}{
            \makecell[{{p{\linewidth}}}]{
                \texttt{\tiny{[GM$|$GM]}}
                {'word': 'ядерне паливо', 'assignment': 'team'}
            }
        }
    }
    & & \\ \\

    \theutterance \stepcounter{utterance}  
    & & & \multicolumn{2}{p{0.3\linewidth}}{
        \cellcolor[rgb]{0.9,0.9,0.9}{
            \makecell[{{p{\linewidth}}}]{
                \texttt{\tiny{[GM$|$GM]}}
                {'word': 'ядерне паливо', 'assignment': 'team'}
            }
        }
    }
    & & \\ \\

    \theutterance \stepcounter{utterance}  
    & & & \multicolumn{2}{p{0.3\linewidth}}{
        \cellcolor[rgb]{0.9,0.9,0.9}{
            \makecell[{{p{\linewidth}}}]{
                \texttt{\tiny{[GM$|$GM]}}
                {'word': 'сонячна енергія', 'assignment': 'opponent'}
            }
        }
    }
    & & \\ \\

    \theutterance \stepcounter{utterance}  
    & & & \multicolumn{2}{p{0.3\linewidth}}{
        \cellcolor[rgb]{0.9,0.9,0.9}{
            \makecell[{{p{\linewidth}}}]{
                \texttt{\tiny{[GM$|$GM]}}
                {'word': 'сонячна енергія', 'assignment': 'opponent'}
            }
        }
    }
    & & \\ \\

    \theutterance \stepcounter{utterance}  
    & & & \multicolumn{2}{p{0.3\linewidth}}{
        \cellcolor[rgb]{0.9,0.9,0.9}{
            \makecell[{{p{\linewidth}}}]{
                \texttt{\tiny{[GM$|$GM]}}
                \texttt{сонячна енергія} \\
            }
        }
    }
    & & \\ \\

    \theutterance \stepcounter{utterance}  
    & & & \multicolumn{2}{p{0.3\linewidth}}{
        \cellcolor[rgb]{0.9,0.9,0.9}{
            \makecell[{{p{\linewidth}}}]{
                \texttt{\tiny{[GM$|$GM]}}
                {'word': 'сом', 'assignment': 'opponent'}
            }
        }
    }
    & & \\ \\

    \theutterance \stepcounter{utterance}  
    & & & \multicolumn{2}{p{0.3\linewidth}}{
        \cellcolor[rgb]{0.9,0.9,0.9}{
            \makecell[{{p{\linewidth}}}]{
                \texttt{\tiny{[GM$|$GM]}}
                {'hidden': {'team': ['цукор', 'лялька', 'квартира', 'корова', 'опера', 'лось', 'неетилований бензин', 'вігвам'], 'innocent': ['білка', 'нирка', 'кориця', 'блюграсс', 'масло', 'кит', 'кубики'], 'opponent': ['фолк', 'кумин', 'шкіра', 'вітер', 'пазл', 'пума'], 'assassin': ['джакс']}, 'revealed': {'team': {'team': ['ядерне паливо'], 'innocent': [], 'opponent': ['сонячна енергія'], 'assassin': []}, 'opponent': {'team': [], 'innocent': [], 'opponent': ['сом'], 'assassin': []}}}
            }
        }
    }
    & & \\ \\

    \theutterance \stepcounter{utterance}  
    & & \multicolumn{4}{p{0.6\linewidth}}{
        \cellcolor[rgb]{0.9,0.9,0.9}{
            \makecell[{{p{\linewidth}}}]{
                \texttt{\tiny{[P1$\langle$GM]}}
                \texttt{Слово ядерне паливо було відгадано правильно. Слово сонячна енергія було відгадано, але це суперницьке слово. Хід вашого товариша по команді закінчився.} \\
\texttt{Тепер надайте ще одну підказку, пов'язану з деякими з ваших командних слів, що залишилися, та список пов'язаних цільових слів. Пам'ятайте: почніть підказку з 'ПІДКАЗКА: ', перейдіть на новий рядок і почніть ваш список цільових слів, розділених комами, з 'ЦІЛІ: '. Зверніть увагу: деякі слова було вилучено зі списків порівняно з попередніми запитами.} \\
\\ 
\texttt{Ваші командні слова, що залишилися: цукор, лялька, квартира, корова, опера, лось, неетилований бензин, вігвам.} \\
\texttt{Слова, що залишилися для вашого суперника: фолк, кумин, шкіра, сом, вітер, пазл, пума.} \\
\texttt{Слова{-}відволікачі, що залишилися: білка, нирка, кориця, блюграсс, масло, кит, кубики.} \\
\texttt{Слова{-}вбивці, що залишилися: джакс.} \\
            }
        }
    }
    & \\ \\

    \theutterance \stepcounter{utterance}  
    & \multicolumn{4}{p{0.6\linewidth}}{
        \cellcolor[rgb]{0.9,0.9,0.9}{
            \makecell[{{p{\linewidth}}}]{
                \texttt{\tiny{[P1$\rangle$GM]}}
                \texttt{ПІДКАЗКА: агропромисловий} \\
\texttt{ЦІЛІ: корова, квартира, вігвам} \\
            }
        }
    }
    & & \\ \\

    \theutterance \stepcounter{utterance}  
    & & & \multicolumn{2}{p{0.3\linewidth}}{
        \cellcolor[rgb]{0.9,0.9,0.9}{
            \makecell[{{p{\linewidth}}}]{
                \texttt{\tiny{[GM$|$GM]}}
                \texttt{агропромисловий} \\
            }
        }
    }
    & & \\ \\

    \theutterance \stepcounter{utterance}  
    & & & \multicolumn{2}{p{0.3\linewidth}}{
        \cellcolor[rgb]{0.9,0.9,0.9}{
            \makecell[{{p{\linewidth}}}]{
                \texttt{\tiny{[GM$|$GM]}}
                ['корова', 'квартира', 'вігвам']
            }
        }
    }
    & & \\ \\

    \theutterance \stepcounter{utterance}  
    & & & \multicolumn{2}{p{0.3\linewidth}}{
        \cellcolor[rgb]{0.9,0.9,0.9}{
            \makecell[{{p{\linewidth}}}]{
                \texttt{\tiny{[GM$|$GM]}}
                {'word': 'корова', 'assignment': 'team'}
            }
        }
    }
    & & \\ \\

    \theutterance \stepcounter{utterance}  
    & & & \multicolumn{2}{p{0.3\linewidth}}{
        \cellcolor[rgb]{0.9,0.9,0.9}{
            \makecell[{{p{\linewidth}}}]{
                \texttt{\tiny{[GM$|$GM]}}
                {'word': 'квартира', 'assignment': 'team'}
            }
        }
    }
    & & \\ \\

    \theutterance \stepcounter{utterance}  
    & & & \multicolumn{2}{p{0.3\linewidth}}{
        \cellcolor[rgb]{0.9,0.9,0.9}{
            \makecell[{{p{\linewidth}}}]{
                \texttt{\tiny{[GM$|$GM]}}
                {'word': 'вігвам', 'assignment': 'team'}
            }
        }
    }
    & & \\ \\

    \theutterance \stepcounter{utterance}  
    & & \multicolumn{4}{p{0.6\linewidth}}{
        \cellcolor[rgb]{0.9,0.9,0.9}{
            \makecell[{{p{\linewidth}}}]{
                \texttt{\tiny{[GM$\rangle$P2]}}
                \texttt{Слово ядерне паливо було відгадано правильно. Слово сонячна енергія було відгадано, але це суперницьке слово. Ваш хід закінчився.} \\
\texttt{Тепер надайте ще один список слів, розділених комами, щонайменше з 1 і щонайбільше з 3 слів із наведеного нижче списку слів, які найкраще пов'язані або найбільш тісно асоціюються зі словом 'агропромисловий'. Пам'ятайте, що потрібно почати вашу відповідь з ВІДГАДКА:. Зауваження: деякі слова було вилучено зі списку порівняно з попередніми запитами.} \\
\\ 
\texttt{нирка, квартира, лось, блюграсс, кориця, кит, білка, джакс, цукор, фолк, пазл, опера, неетилований бензин, лялька, вітер, кумин, корова, вігвам, шкіра, масло, пума, кубики} \\
            }
        }
    }
    & \\ \\

    \theutterance \stepcounter{utterance}  
    & & & \multicolumn{4}{p{0.6\linewidth}}{
        \cellcolor[rgb]{0.9,0.9,0.9}{
            \makecell[{{p{\linewidth}}}]{
                \texttt{\tiny{[GM$\langle$P2]}}
                \texttt{ВІДГАДКА: корова, кумин, масло} \\
            }
        }
    }
     \\ \\

    \theutterance \stepcounter{utterance}  
    & & & \multicolumn{2}{p{0.3\linewidth}}{
        \cellcolor[rgb]{0.9,0.9,0.9}{
            \makecell[{{p{\linewidth}}}]{
                \texttt{\tiny{[GM$|$GM]}}
                ['корова', 'кумин', 'масло']
            }
        }
    }
    & & \\ \\

    \theutterance \stepcounter{utterance}  
    & & & \multicolumn{2}{p{0.3\linewidth}}{
        \cellcolor[rgb]{0.9,0.9,0.9}{
            \makecell[{{p{\linewidth}}}]{
                \texttt{\tiny{[GM$|$GM]}}
                {'word': 'корова', 'assignment': 'team'}
            }
        }
    }
    & & \\ \\

    \theutterance \stepcounter{utterance}  
    & & & \multicolumn{2}{p{0.3\linewidth}}{
        \cellcolor[rgb]{0.9,0.9,0.9}{
            \makecell[{{p{\linewidth}}}]{
                \texttt{\tiny{[GM$|$GM]}}
                {'word': 'корова', 'assignment': 'team'}
            }
        }
    }
    & & \\ \\

    \theutterance \stepcounter{utterance}  
    & & & \multicolumn{2}{p{0.3\linewidth}}{
        \cellcolor[rgb]{0.9,0.9,0.9}{
            \makecell[{{p{\linewidth}}}]{
                \texttt{\tiny{[GM$|$GM]}}
                {'word': 'кумин', 'assignment': 'opponent'}
            }
        }
    }
    & & \\ \\

    \theutterance \stepcounter{utterance}  
    & & & \multicolumn{2}{p{0.3\linewidth}}{
        \cellcolor[rgb]{0.9,0.9,0.9}{
            \makecell[{{p{\linewidth}}}]{
                \texttt{\tiny{[GM$|$GM]}}
                \texttt{кумин} \\
            }
        }
    }
    & & \\ \\

    \theutterance \stepcounter{utterance}  
    & & & \multicolumn{2}{p{0.3\linewidth}}{
        \cellcolor[rgb]{0.9,0.9,0.9}{
            \makecell[{{p{\linewidth}}}]{
                \texttt{\tiny{[GM$|$GM]}}
                {'word': 'пума', 'assignment': 'opponent'}
            }
        }
    }
    & & \\ \\

    \theutterance \stepcounter{utterance}  
    & & & \multicolumn{2}{p{0.3\linewidth}}{
        \cellcolor[rgb]{0.9,0.9,0.9}{
            \makecell[{{p{\linewidth}}}]{
                \texttt{\tiny{[GM$|$GM]}}
                {'hidden': {'team': ['цукор', 'лялька', 'квартира', 'опера', 'лось', 'неетилований бензин', 'вігвам'], 'innocent': ['білка', 'нирка', 'кориця', 'блюграсс', 'масло', 'кит', 'кубики'], 'opponent': ['фолк', 'шкіра', 'вітер', 'пазл'], 'assassin': ['джакс']}, 'revealed': {'team': {'team': ['ядерне паливо', 'корова'], 'innocent': [], 'opponent': ['сонячна енергія', 'кумин'], 'assassin': []}, 'opponent': {'team': [], 'innocent': [], 'opponent': ['сом', 'пума'], 'assassin': []}}}
            }
        }
    }
    & & \\ \\

    \theutterance \stepcounter{utterance}  
    & & \multicolumn{4}{p{0.6\linewidth}}{
        \cellcolor[rgb]{0.9,0.9,0.9}{
            \makecell[{{p{\linewidth}}}]{
                \texttt{\tiny{[P1$\langle$GM]}}
                \texttt{Слово корова було відгадано правильно. Слово кумин було відгадано, але це суперницьке слово. Хід вашого товариша по команді закінчився.} \\
\texttt{Тепер надайте ще одну підказку, пов'язану з деякими з ваших командних слів, що залишилися, та список пов'язаних цільових слів. Пам'ятайте: почніть підказку з 'ПІДКАЗКА: ', перейдіть на новий рядок і почніть ваш список цільових слів, розділених комами, з 'ЦІЛІ: '. Зверніть увагу: деякі слова було вилучено зі списків порівняно з попередніми запитами.} \\
\\ 
\texttt{Ваші командні слова, що залишилися: цукор, лялька, квартира, опера, лось, неетилований бензин, вігвам.} \\
\texttt{Слова, що залишилися для вашого суперника: фолк, шкіра, вітер, пазл, пума.} \\
\texttt{Слова{-}відволікачі, що залишилися: білка, нирка, кориця, блюграсс, масло, кит, кубики.} \\
\texttt{Слова{-}вбивці, що залишилися: джакс.} \\
            }
        }
    }
    & \\ \\

    \theutterance \stepcounter{utterance}  
    & \multicolumn{4}{p{0.6\linewidth}}{
        \cellcolor[rgb]{0.9,0.9,0.9}{
            \makecell[{{p{\linewidth}}}]{
                \texttt{\tiny{[P1$\rangle$GM]}}
                \texttt{ПІДКАЗКА: культурний} \\
\texttt{ЦІЛІ: опера, фолк, пазл} \\
            }
        }
    }
    & & \\ \\

    \theutterance \stepcounter{utterance}  
    & & & \multicolumn{2}{p{0.3\linewidth}}{
        \cellcolor[rgb]{0.9,0.9,0.9}{
            \makecell[{{p{\linewidth}}}]{
                \texttt{\tiny{[GM$|$GM]}}
                \texttt{культурний} \\
            }
        }
    }
    & & \\ \\

    \theutterance \stepcounter{utterance}  
    & & & \multicolumn{2}{p{0.3\linewidth}}{
        \cellcolor[rgb]{0.9,0.9,0.9}{
            \makecell[{{p{\linewidth}}}]{
                \texttt{\tiny{[GM$|$GM]}}
                ['опера', 'фолк', 'пазл']
            }
        }
    }
    & & \\ \\

    \theutterance \stepcounter{utterance}  
    & & & \multicolumn{2}{p{0.3\linewidth}}{
        \cellcolor[rgb]{0.9,0.9,0.9}{
            \makecell[{{p{\linewidth}}}]{
                \texttt{\tiny{[GM$|$GM]}}
                {'word': 'опера', 'assignment': 'team'}
            }
        }
    }
    & & \\ \\

    \theutterance \stepcounter{utterance}  
    & & & \multicolumn{2}{p{0.3\linewidth}}{
        \cellcolor[rgb]{0.9,0.9,0.9}{
            \makecell[{{p{\linewidth}}}]{
                \texttt{\tiny{[GM$|$GM]}}
                {'word': 'фолк', 'assignment': 'opponent'}
            }
        }
    }
    & & \\ \\

    \theutterance \stepcounter{utterance}  
    & & & \multicolumn{2}{p{0.3\linewidth}}{
        \cellcolor[rgb]{0.9,0.9,0.9}{
            \makecell[{{p{\linewidth}}}]{
                \texttt{\tiny{[GM$|$GM]}}
                {'word': 'пазл', 'assignment': 'opponent'}
            }
        }
    }
    & & \\ \\

    \theutterance \stepcounter{utterance}  
    & & \multicolumn{4}{p{0.6\linewidth}}{
        \cellcolor[rgb]{0.9,0.9,0.9}{
            \makecell[{{p{\linewidth}}}]{
                \texttt{\tiny{[GM$\rangle$P2]}}
                \texttt{Слово корова було відгадано правильно. Слово кумин було відгадано, але це суперницьке слово. Ваш хід закінчився.} \\
\texttt{Тепер надайте ще один список слів, розділених комами, щонайменше з 1 і щонайбільше з 3 слів із наведеного нижче списку слів, які найкраще пов'язані або найбільш тісно асоціюються зі словом 'культурний'. Пам'ятайте, що потрібно почати вашу відповідь з ВІДГАДКА:. Зауваження: деякі слова було вилучено зі списку порівняно з попередніми запитами.} \\
\\ 
\texttt{нирка, квартира, лось, блюграсс, кориця, кит, білка, джакс, цукор, фолк, пазл, опера, неетилований бензин, лялька, вітер, вігвам, шкіра, масло, кубики} \\
            }
        }
    }
    & \\ \\

    \theutterance \stepcounter{utterance}  
    & & & \multicolumn{4}{p{0.6\linewidth}}{
        \cellcolor[rgb]{0.9,0.9,0.9}{
            \makecell[{{p{\linewidth}}}]{
                \texttt{\tiny{[GM$\langle$P2]}}
                \texttt{ВІДГАДКА: фолк, опера, пазл} \\
            }
        }
    }
     \\ \\

    \theutterance \stepcounter{utterance}  
    & & & \multicolumn{2}{p{0.3\linewidth}}{
        \cellcolor[rgb]{0.9,0.9,0.9}{
            \makecell[{{p{\linewidth}}}]{
                \texttt{\tiny{[GM$|$GM]}}
                ['фолк', 'опера', 'пазл']
            }
        }
    }
    & & \\ \\

    \theutterance \stepcounter{utterance}  
    & & & \multicolumn{2}{p{0.3\linewidth}}{
        \cellcolor[rgb]{0.9,0.9,0.9}{
            \makecell[{{p{\linewidth}}}]{
                \texttt{\tiny{[GM$|$GM]}}
                {'word': 'фолк', 'assignment': 'opponent'}
            }
        }
    }
    & & \\ \\

    \theutterance \stepcounter{utterance}  
    & & & \multicolumn{2}{p{0.3\linewidth}}{
        \cellcolor[rgb]{0.9,0.9,0.9}{
            \makecell[{{p{\linewidth}}}]{
                \texttt{\tiny{[GM$|$GM]}}
                {'word': 'фолк', 'assignment': 'opponent'}
            }
        }
    }
    & & \\ \\

    \theutterance \stepcounter{utterance}  
    & & & \multicolumn{2}{p{0.3\linewidth}}{
        \cellcolor[rgb]{0.9,0.9,0.9}{
            \makecell[{{p{\linewidth}}}]{
                \texttt{\tiny{[GM$|$GM]}}
                \texttt{фолк} \\
            }
        }
    }
    & & \\ \\

    \theutterance \stepcounter{utterance}  
    & & & \multicolumn{2}{p{0.3\linewidth}}{
        \cellcolor[rgb]{0.9,0.9,0.9}{
            \makecell[{{p{\linewidth}}}]{
                \texttt{\tiny{[GM$|$GM]}}
                {'word': 'вітер', 'assignment': 'opponent'}
            }
        }
    }
    & & \\ \\

    \theutterance \stepcounter{utterance}  
    & & & \multicolumn{2}{p{0.3\linewidth}}{
        \cellcolor[rgb]{0.9,0.9,0.9}{
            \makecell[{{p{\linewidth}}}]{
                \texttt{\tiny{[GM$|$GM]}}
                {'hidden': {'team': ['цукор', 'лялька', 'квартира', 'опера', 'лось', 'неетилований бензин', 'вігвам'], 'innocent': ['білка', 'нирка', 'кориця', 'блюграсс', 'масло', 'кит', 'кубики'], 'opponent': ['шкіра', 'пазл'], 'assassin': ['джакс']}, 'revealed': {'team': {'team': ['ядерне паливо', 'корова'], 'innocent': [], 'opponent': ['сонячна енергія', 'кумин', 'фолк'], 'assassin': []}, 'opponent': {'team': [], 'innocent': [], 'opponent': ['сом', 'пума', 'вітер'], 'assassin': []}}}
            }
        }
    }
    & & \\ \\

    \theutterance \stepcounter{utterance}  
    & & \multicolumn{4}{p{0.6\linewidth}}{
        \cellcolor[rgb]{0.9,0.9,0.9}{
            \makecell[{{p{\linewidth}}}]{
                \texttt{\tiny{[P1$\langle$GM]}}
                \texttt{Слово фолк було відгадано, але це суперницьке слово. Хід вашого товариша по команді закінчився.} \\
\texttt{Тепер надайте ще одну підказку, пов'язану з деякими з ваших командних слів, що залишилися, та список пов'язаних цільових слів. Пам'ятайте: почніть підказку з 'ПІДКАЗКА: ', перейдіть на новий рядок і почніть ваш список цільових слів, розділених комами, з 'ЦІЛІ: '. Зверніть увагу: деякі слова було вилучено зі списків порівняно з попередніми запитами.} \\
\\ 
\texttt{Ваші командні слова, що залишилися: цукор, лялька, квартира, опера, лось, неетилований бензин, вігвам.} \\
\texttt{Слова, що залишилися для вашого суперника: шкіра, вітер, пазл.} \\
\texttt{Слова{-}відволікачі, що залишилися: білка, нирка, кориця, блюграсс, масло, кит, кубики.} \\
\texttt{Слова{-}вбивці, що залишилися: джакс.} \\
            }
        }
    }
    & \\ \\

    \theutterance \stepcounter{utterance}  
    & \multicolumn{4}{p{0.6\linewidth}}{
        \cellcolor[rgb]{0.9,0.9,0.9}{
            \makecell[{{p{\linewidth}}}]{
                \texttt{\tiny{[P1$\rangle$GM]}}
                \texttt{ПІДКАЗКА: житло} \\
\texttt{ЦІЛІ: квартира, вігвам} \\
            }
        }
    }
    & & \\ \\

    \theutterance \stepcounter{utterance}  
    & & & \multicolumn{2}{p{0.3\linewidth}}{
        \cellcolor[rgb]{0.9,0.9,0.9}{
            \makecell[{{p{\linewidth}}}]{
                \texttt{\tiny{[GM$|$GM]}}
                \texttt{житло} \\
            }
        }
    }
    & & \\ \\

    \theutterance \stepcounter{utterance}  
    & & & \multicolumn{2}{p{0.3\linewidth}}{
        \cellcolor[rgb]{0.9,0.9,0.9}{
            \makecell[{{p{\linewidth}}}]{
                \texttt{\tiny{[GM$|$GM]}}
                ['квартира', 'вігвам']
            }
        }
    }
    & & \\ \\

    \theutterance \stepcounter{utterance}  
    & & & \multicolumn{2}{p{0.3\linewidth}}{
        \cellcolor[rgb]{0.9,0.9,0.9}{
            \makecell[{{p{\linewidth}}}]{
                \texttt{\tiny{[GM$|$GM]}}
                {'word': 'квартира', 'assignment': 'team'}
            }
        }
    }
    & & \\ \\

    \theutterance \stepcounter{utterance}  
    & & & \multicolumn{2}{p{0.3\linewidth}}{
        \cellcolor[rgb]{0.9,0.9,0.9}{
            \makecell[{{p{\linewidth}}}]{
                \texttt{\tiny{[GM$|$GM]}}
                {'word': 'вігвам', 'assignment': 'team'}
            }
        }
    }
    & & \\ \\

    \theutterance \stepcounter{utterance}  
    & & \multicolumn{4}{p{0.6\linewidth}}{
        \cellcolor[rgb]{0.9,0.9,0.9}{
            \makecell[{{p{\linewidth}}}]{
                \texttt{\tiny{[GM$\rangle$P2]}}
                \texttt{Слово фолк було відгадано, але це суперницьке слово. Ваш хід закінчився.} \\
\texttt{Тепер надайте ще один список слів, розділених комами, щонайменше з 1 і щонайбільше з 2 слів із наведеного нижче списку слів, які найкраще пов'язані або найбільш тісно асоціюються зі словом 'житло'. Пам'ятайте, що потрібно почати вашу відповідь з ВІДГАДКА:. Зауваження: деякі слова було вилучено зі списку порівняно з попередніми запитами.} \\
\\ 
\texttt{нирка, квартира, лось, блюграсс, кориця, кит, білка, джакс, цукор, пазл, опера, неетилований бензин, лялька, вігвам, шкіра, масло, кубики} \\
            }
        }
    }
    & \\ \\

    \theutterance \stepcounter{utterance}  
    & & & \multicolumn{4}{p{0.6\linewidth}}{
        \cellcolor[rgb]{0.9,0.9,0.9}{
            \makecell[{{p{\linewidth}}}]{
                \texttt{\tiny{[GM$\langle$P2]}}
                \texttt{ВІДГАДКА: квартира, вігвам} \\
            }
        }
    }
     \\ \\

    \theutterance \stepcounter{utterance}  
    & & & \multicolumn{2}{p{0.3\linewidth}}{
        \cellcolor[rgb]{0.9,0.9,0.9}{
            \makecell[{{p{\linewidth}}}]{
                \texttt{\tiny{[GM$|$GM]}}
                ['квартира', 'вігвам']
            }
        }
    }
    & & \\ \\

    \theutterance \stepcounter{utterance}  
    & & & \multicolumn{2}{p{0.3\linewidth}}{
        \cellcolor[rgb]{0.9,0.9,0.9}{
            \makecell[{{p{\linewidth}}}]{
                \texttt{\tiny{[GM$|$GM]}}
                {'word': 'квартира', 'assignment': 'team'}
            }
        }
    }
    & & \\ \\

    \theutterance \stepcounter{utterance}  
    & & & \multicolumn{2}{p{0.3\linewidth}}{
        \cellcolor[rgb]{0.9,0.9,0.9}{
            \makecell[{{p{\linewidth}}}]{
                \texttt{\tiny{[GM$|$GM]}}
                {'word': 'квартира', 'assignment': 'team'}
            }
        }
    }
    & & \\ \\

    \theutterance \stepcounter{utterance}  
    & & & \multicolumn{2}{p{0.3\linewidth}}{
        \cellcolor[rgb]{0.9,0.9,0.9}{
            \makecell[{{p{\linewidth}}}]{
                \texttt{\tiny{[GM$|$GM]}}
                {'word': 'вігвам', 'assignment': 'team'}
            }
        }
    }
    & & \\ \\

    \theutterance \stepcounter{utterance}  
    & & & \multicolumn{2}{p{0.3\linewidth}}{
        \cellcolor[rgb]{0.9,0.9,0.9}{
            \makecell[{{p{\linewidth}}}]{
                \texttt{\tiny{[GM$|$GM]}}
                {'word': 'вігвам', 'assignment': 'team'}
            }
        }
    }
    & & \\ \\

    \theutterance \stepcounter{utterance}  
    & & & \multicolumn{2}{p{0.3\linewidth}}{
        \cellcolor[rgb]{0.9,0.9,0.9}{
            \makecell[{{p{\linewidth}}}]{
                \texttt{\tiny{[GM$|$GM]}}
                {'word': 'шкіра', 'assignment': 'opponent'}
            }
        }
    }
    & & \\ \\

    \theutterance \stepcounter{utterance}  
    & & & \multicolumn{2}{p{0.3\linewidth}}{
        \cellcolor[rgb]{0.9,0.9,0.9}{
            \makecell[{{p{\linewidth}}}]{
                \texttt{\tiny{[GM$|$GM]}}
                {'hidden': {'team': ['цукор', 'лялька', 'опера', 'лось', 'неетилований бензин'], 'innocent': ['білка', 'нирка', 'кориця', 'блюграсс', 'масло', 'кит', 'кубики'], 'opponent': ['пазл'], 'assassin': ['джакс']}, 'revealed': {'team': {'team': ['ядерне паливо', 'корова', 'квартира', 'вігвам'], 'innocent': [], 'opponent': ['сонячна енергія', 'кумин', 'фолк'], 'assassin': []}, 'opponent': {'team': [], 'innocent': [], 'opponent': ['сом', 'пума', 'вітер', 'шкіра'], 'assassin': []}}}
            }
        }
    }
    & & \\ \\

    \theutterance \stepcounter{utterance}  
    & & \multicolumn{4}{p{0.6\linewidth}}{
        \cellcolor[rgb]{0.9,0.9,0.9}{
            \makecell[{{p{\linewidth}}}]{
                \texttt{\tiny{[P1$\langle$GM]}}
                \texttt{Слова квартира, вігвам було відгадано правильно. Хід вашого товариша по команді закінчився.} \\
\texttt{Тепер надайте ще одну підказку, пов'язану з деякими з ваших командних слів, що залишилися, та список пов'язаних цільових слів. Пам'ятайте: почніть підказку з 'ПІДКАЗКА: ', перейдіть на новий рядок і почніть ваш список цільових слів, розділених комами, з 'ЦІЛІ: '. Зверніть увагу: деякі слова було вилучено зі списків порівняно з попередніми запитами.} \\
\\ 
\texttt{Ваші командні слова, що залишилися: цукор, лялька, опера, лось, неетилований бензин.} \\
\texttt{Слова, що залишилися для вашого суперника: шкіра, пазл.} \\
\texttt{Слова{-}відволікачі, що залишилися: білка, нирка, кориця, блюграсс, масло, кит, кубики.} \\
\texttt{Слова{-}вбивці, що залишилися: джакс.} \\
            }
        }
    }
    & \\ \\

    \theutterance \stepcounter{utterance}  
    & \multicolumn{4}{p{0.6\linewidth}}{
        \cellcolor[rgb]{0.9,0.9,0.9}{
            \makecell[{{p{\linewidth}}}]{
                \texttt{\tiny{[P1$\rangle$GM]}}
                \texttt{ПІДКАЗКА: мистецтво} \\
\texttt{ЦІЛІ: опера} \\
            }
        }
    }
    & & \\ \\

    \theutterance \stepcounter{utterance}  
    & & & \multicolumn{2}{p{0.3\linewidth}}{
        \cellcolor[rgb]{0.9,0.9,0.9}{
            \makecell[{{p{\linewidth}}}]{
                \texttt{\tiny{[GM$|$GM]}}
                \texttt{мистецтво} \\
            }
        }
    }
    & & \\ \\

    \theutterance \stepcounter{utterance}  
    & & & \multicolumn{2}{p{0.3\linewidth}}{
        \cellcolor[rgb]{0.9,0.9,0.9}{
            \makecell[{{p{\linewidth}}}]{
                \texttt{\tiny{[GM$|$GM]}}
                ['опера']
            }
        }
    }
    & & \\ \\

    \theutterance \stepcounter{utterance}  
    & & & \multicolumn{2}{p{0.3\linewidth}}{
        \cellcolor[rgb]{0.9,0.9,0.9}{
            \makecell[{{p{\linewidth}}}]{
                \texttt{\tiny{[GM$|$GM]}}
                {'word': 'опера', 'assignment': 'team'}
            }
        }
    }
    & & \\ \\

    \theutterance \stepcounter{utterance}  
    & & \multicolumn{4}{p{0.6\linewidth}}{
        \cellcolor[rgb]{0.9,0.9,0.9}{
            \makecell[{{p{\linewidth}}}]{
                \texttt{\tiny{[GM$\rangle$P2]}}
                \texttt{Слова квартира, вігвам було відгадано правильно. Ваш хід закінчився.} \\
\texttt{Тепер надайте ще один список слів, розділених комами, щонайменше з 1 і щонайбільше з 1 слів із наведеного нижче списку слів, які найкраще пов'язані або найбільш тісно асоціюються зі словом 'мистецтво'. Пам'ятайте, що потрібно почати вашу відповідь з ВІДГАДКА:. Зауваження: деякі слова було вилучено зі списку порівняно з попередніми запитами.} \\
\\ 
\texttt{нирка, лось, блюграсс, кориця, кит, білка, джакс, цукор, пазл, опера, неетилований бензин, лялька, масло, кубики} \\
            }
        }
    }
    & \\ \\

    \theutterance \stepcounter{utterance}  
    & & & \multicolumn{4}{p{0.6\linewidth}}{
        \cellcolor[rgb]{0.9,0.9,0.9}{
            \makecell[{{p{\linewidth}}}]{
                \texttt{\tiny{[GM$\langle$P2]}}
                \texttt{ВІДГАДКА: опера} \\
            }
        }
    }
     \\ \\

    \theutterance \stepcounter{utterance}  
    & & & \multicolumn{2}{p{0.3\linewidth}}{
        \cellcolor[rgb]{0.9,0.9,0.9}{
            \makecell[{{p{\linewidth}}}]{
                \texttt{\tiny{[GM$|$GM]}}
                ['опера']
            }
        }
    }
    & & \\ \\

    \theutterance \stepcounter{utterance}  
    & & & \multicolumn{2}{p{0.3\linewidth}}{
        \cellcolor[rgb]{0.9,0.9,0.9}{
            \makecell[{{p{\linewidth}}}]{
                \texttt{\tiny{[GM$|$GM]}}
                {'word': 'опера', 'assignment': 'team'}
            }
        }
    }
    & & \\ \\

    \theutterance \stepcounter{utterance}  
    & & & \multicolumn{2}{p{0.3\linewidth}}{
        \cellcolor[rgb]{0.9,0.9,0.9}{
            \makecell[{{p{\linewidth}}}]{
                \texttt{\tiny{[GM$|$GM]}}
                {'word': 'опера', 'assignment': 'team'}
            }
        }
    }
    & & \\ \\

    \theutterance \stepcounter{utterance}  
    & & & \multicolumn{2}{p{0.3\linewidth}}{
        \cellcolor[rgb]{0.9,0.9,0.9}{
            \makecell[{{p{\linewidth}}}]{
                \texttt{\tiny{[GM$|$GM]}}
                {'word': 'пазл', 'assignment': 'opponent'}
            }
        }
    }
    & & \\ \\

    \theutterance \stepcounter{utterance}  
    & & & \multicolumn{2}{p{0.3\linewidth}}{
        \cellcolor[rgb]{0.9,0.9,0.9}{
            \makecell[{{p{\linewidth}}}]{
                \texttt{\tiny{[GM$|$GM]}}
                {'hidden': {'team': ['цукор', 'лялька', 'лось', 'неетилований бензин'], 'innocent': ['білка', 'нирка', 'кориця', 'блюграсс', 'масло', 'кит', 'кубики'], 'opponent': [], 'assassin': ['джакс']}, 'revealed': {'team': {'team': ['ядерне паливо', 'корова', 'квартира', 'вігвам', 'опера'], 'innocent': [], 'opponent': ['сонячна енергія', 'кумин', 'фолк'], 'assassin': []}, 'opponent': {'team': [], 'innocent': [], 'opponent': ['сом', 'пума', 'вітер', 'шкіра', 'пазл'], 'assassin': []}}}
            }
        }
    }
    & & \\ \\

    \theutterance \stepcounter{utterance}  
    & & \multicolumn{4}{p{0.6\linewidth}}{
        \cellcolor[rgb]{0.9,0.9,0.9}{
            \makecell[{{p{\linewidth}}}]{
                \texttt{\tiny{[P1$\langle$GM]}}
                \texttt{Слово опера було відгадано правильно. Хід вашого товариша по команді закінчився.} \\
\texttt{Тепер надайте ще одну підказку, пов'язану з деякими з ваших командних слів, що залишилися, та список пов'язаних цільових слів. Пам'ятайте: почніть підказку з 'ПІДКАЗКА: ', перейдіть на новий рядок і почніть ваш список цільових слів, розділених комами, з 'ЦІЛІ: '. Зверніть увагу: деякі слова було вилучено зі списків порівняно з попередніми запитами.} \\
\\ 
\texttt{Ваші командні слова, що залишилися: цукор, лялька, лось, неетилований бензин.} \\
\texttt{Слова, що залишилися для вашого суперника: пазл.} \\
\texttt{Слова{-}відволікачі, що залишилися: білка, нирка, кориця, блюграсс, масло, кит, кубики.} \\
\texttt{Слова{-}вбивці, що залишилися: джакс.} \\
            }
        }
    }
    & \\ \\

    \theutterance \stepcounter{utterance}  
    & \multicolumn{4}{p{0.6\linewidth}}{
        \cellcolor[rgb]{0.9,0.9,0.9}{
            \makecell[{{p{\linewidth}}}]{
                \texttt{\tiny{[P1$\rangle$GM]}}
                \texttt{ПІДКАЗКА: транспорт} \\
\texttt{ЦІЛІ: неетилований бензин} \\
            }
        }
    }
    & & \\ \\

    \theutterance \stepcounter{utterance}  
    & & & \multicolumn{2}{p{0.3\linewidth}}{
        \cellcolor[rgb]{0.9,0.9,0.9}{
            \makecell[{{p{\linewidth}}}]{
                \texttt{\tiny{[GM$|$GM]}}
                \texttt{транспорт} \\
            }
        }
    }
    & & \\ \\

    \theutterance \stepcounter{utterance}  
    & & & \multicolumn{2}{p{0.3\linewidth}}{
        \cellcolor[rgb]{0.9,0.9,0.9}{
            \makecell[{{p{\linewidth}}}]{
                \texttt{\tiny{[GM$|$GM]}}
                ['неетилований бензин']
            }
        }
    }
    & & \\ \\

    \theutterance \stepcounter{utterance}  
    & & & \multicolumn{2}{p{0.3\linewidth}}{
        \cellcolor[rgb]{0.9,0.9,0.9}{
            \makecell[{{p{\linewidth}}}]{
                \texttt{\tiny{[GM$|$GM]}}
                {'word': 'неетилований бензин', 'assignment': 'team'}
            }
        }
    }
    & & \\ \\

    \theutterance \stepcounter{utterance}  
    & & & \multicolumn{2}{p{0.3\linewidth}}{
        \cellcolor[rgb]{0.9,0.9,0.9}{
            \makecell[{{p{\linewidth}}}]{
                \texttt{\tiny{[GM$|$GM]}}
                \texttt{opponent has won} \\
            }
        }
    }
    & & \\ \\

\end{supertabular}
}

\end{document}
